% Chapter 3
\chapter{Technologies and Frameworks} % Main chapter title
\label{Chapter3}

%----------------------------------------------------------------------------------------

The IIITA Timetabling System utilizes a modern, browser-based stack selected to prioritize high-performance data processing and strict referential integrity. This chapter outlines the technical foundation, focusing on tools that facilitate complex client-side computation and the normalization of institutional data.

\begin{table}[htbp]
\centering
\caption{Summary of the Technology Stack}
\label{tab:tech_stack}
\renewcommand{\arraystretch}{1.3}
\begin{tabularx}{\textwidth}{@{} l l X @{}}
\toprule
\textbf{Layer} & \textbf{Technology} & \textbf{Primary Function in TTGen} \\
\midrule
Frontend Framework & React 19 & Rendering dense, interactive 5-day time grids. \\
Programming Language & TypeScript & Enforcing strict data contracts for slots and genes. \\
Database & PostgreSQL & Maintaining integrity across rooms, faculty, and subjects. \\
Backend / BaaS & Supabase & API gateway and secure database interfacing. \\
Excel Parsing & SheetJS (xlsx) & Heuristic extraction of semantic meaning from legacy files. \\
\bottomrule
\end{tabularx}
\end{table}

%-----------------------------------
\section{Frontend Architecture: React 19 and TypeScript}

\begin{figure}[htbp]
\centering
\begin{tikzpicture}[
    node distance=1.5cm and 2.5cm,
    client/.style={rectangle, draw, fill=blue!10, rounded corners, text width=3cm, align=center, minimum height=1.5cm, font=\sffamily\small},
    api/.style={rectangle, draw, fill=green!10, text width=3.5cm, align=center, minimum height=1.5cm, font=\sffamily\small},
    db/.style={cylinder, draw, shape border rotate=90, aspect=0.25, fill=orange!10, text width=2.5cm, align=center, minimum height=2cm, font=\sffamily\small},
    arrow/.style={draw, thick, <->, >=stealth}
]

% Nodes
\node[client] (react) {\textbf{Client Application}\\React 19 / Vite\\(Role-Based UI)};
\node[api, right=of react] (supabase) {\textbf{Supabase API Gateway}\\PostgREST\\Auth Service};
\node[db, right=of supabase] (postgres) {\textbf{Data Storage}\\PostgreSQL\\(Relational Schema)};

% Arrows
\draw[arrow] (react) -- (supabase) node[midway, above] {\scriptsize REST/JSON};
\draw[arrow] (supabase) -- (postgres) node[midway, above] {\scriptsize SQL Queries};

\end{tikzpicture}
\caption{High-Level System Architecture and Component Interaction}
\label{fig:system_architecture}
\end{figure}

The user interface is created with the ability to handle dense data interactions without a drop in performance. This is accomplished through the use of React's virtual DOM and TypeScript static type checking.

\subsection{Performance and Data Reliability}

For expensive operations like $O(N^2)$ clash detection, memoization of these logic-driven functions is done using the \texttt{useMemo} hook, such that these functions are run only if the data from the timetable has been modified. TypeScript is utilized for strict data interfaces, like \texttt{ParsedSlot}.

\begin{verbatim}
// Example TypeScript Interface for Solver Data Contracts
interface ParsedSlot {
    id: string;
    day: string;
    startTime: string;
    subjectCode: string;
    type: 'Lecture' | 'Tutorial' | 'Practical';
    roomName: string;
    section: string;
}
\end{verbatim}

%-----------------------------------
\section{Construction and Styling: Vite and Tailwind CSS}

Both the construction and styling processes are designed to facilitate fast iterations and immediate feedback loops.
\begin{itemize}
    \item \textbf{Vite:} Utilizes native ES Modules, enabling hot module replacement (HMR) almost instantaneously, which is essential during the optimization of grid generation algorithms.
    \item \textbf{Tailwind CSS:} The utility-based CSS framework makes it possible for the application to provide real-time styling feedback, such as a red border around conflicting sessions.
\end{itemize}

%-----------------------------------
\section{Back-end and Data Ingestion: Supabase}

Static files are replaced with a central PostgreSQL database provided by Supabase with referential integrity enforced using foreign keys.
\begin{itemize}
    \item \textbf{Heuristic Parsing:} The \texttt{xlsx} (SheetJS) package is utilized to parse institutional spreadsheets in their raw form. This information is interpreted using an in-house heuristic engine to derive its semantics, such as credit structures in the form of L-T-P and faculty mapping.
    \item \textbf{Authentication:} Authentication for administrators is achieved through Supabase Auth, which is implemented with state management through a global \texttt{AuthContext}.
\end{itemize}

%-----------------------------------
\section{Solver Engine: Architecture of Pure TypeScript}

The most important aspect of the current solution is its pure-TypeScript solver engine, which works directly in the browser. It doesn’t require any optimizations on the server side, enabling live generation of schedules. Solver consists of several modules:

\begin{itemize}
    \item \textbf{\texttt{dataPrep.ts}}: Breaks up high-level credits into individual sessions, applying the 2+1 rule of lectures.
    \item \textbf{\texttt{constraints.ts}}: Checks the solutions based on ten hard and three soft constraints.
    \item \textbf{\texttt{mutations.ts}}: Contains five different mutation operators, such as session day changing and Gaussian time shift.
    \item \textbf{\texttt{solver.ts}}: Works as a control center, managing the (1+1)-ES process, elitism, and restarts from stagnation.
    \item \textbf{\texttt{localSearch.ts}}: Applies local search operations to refine the schedule and resolve specific conflicts.
\end{itemize}

%-----------------------------------
\section{External Integration: Reporting and Exporting}

For workflow efficiency purposes, the system has the following integrations with external reporting systems:
\begin{itemize}
    \item \textbf{Google Sheets API:} Enables exportation of schedules to Google Drive in a frozen header format with session coloring.
    \item \textbf{PDF Output:} Makes use of \texttt{html-to-image} and \texttt{jsPDF} libraries for snapshot output of the timetable into PDF form.
\end{itemize}